\section{Fourier action and integrated constraints}

\stepcounter{equation}\begin{equation}\tag{WD-01}\label{eq:WD-01}
\Gamma^{i(1)}{}_{jk}=\tfrac12(\partial_jh_{ik}+\partial_kh_{ij}-\partial_ih_{jk}),\quad R^{[1]}=\partial_i\partial_jh_{ij}-\partial^2h_{ii}.
\end{equation}

The second coefficient of the integrated $\sqrt g R$ Gamma--Gamma form gives the spatial term of (W-01). The first lapse perturbation multiplies $R^{[1]}$. With the wavevector along $z$, $R^{[1]}=k^2(h_{11}+h_{22})$; the shift enters through $v$, retaining all constraints. The momentum equations are proportional to $p_{13}=p_{23}=p_{33}=0$ and the linear Hamiltonian condition to $h_{11}+h_{22}=0$. These four secondaries and four primaries mutually Poisson commute.

For $k\ne0$, quotienting gauge directions and taking $h_{11}=-h_{22}=A$, $h_{12}=B$ gives (W-02). For any nonzero lattice wavevector the same algebra holds in an orthonormal polarization basis. The basis rotation need not be a lattice automorphism: the original wavevector and polarization tensors are restored.

\stepcounter{equation}\begin{equation}\tag{WD-02}\label{eq:WD-02}
\begin{aligned}R+K^2-K_{ij}K^{ij}&=0,\qquad D_j(K^{ij}-g^{ij}K)=0,\\ K^{[1]}&=hI+\Sigma+\tfrac12\dot\gamma,\qquad\left\langle R^{[2]}\right\rangle=-\tfrac1{4q_0^2}\left\langle(\partial_a\gamma_{ij})^2\right\rangle.\end{aligned}
\end{equation}

Zero-mean TT fields have vanishing averaged cross terms with constant $h,\Sigma$. The linear curvature of the second-order metric correction integrates to zero as a divergence, and because the background $K$ vanishes, the second-order correction to $K$ does not enter the quadratic constraint linearly. The averaged Hamiltonian constraint is therefore $6h^2-\mathrm{tr}\Sigma^2-\langle\dot\gamma^2\rangle/4-\langle(\partial\gamma)^2\rangle/(4q_0^2)=0$.

\stepcounter{equation}\begin{equation}\tag{WD-03}\label{eq:WD-03}
\mathcal J[\xi]=\int_\Sigma\pi^{ij}\mathcal L_\xi g_{ij},\quad \xi=\partial_a:\qquad \mathcal J^{[2]}[\partial_a]\ \propto\ \left\langle\dot\gamma_{ij}\partial_a\gamma_{ij}\right\rangle.
\end{equation}

Because background $\pi=0$ and $g$ is constant, linear terms in the second-order fields vanish, leaving (WD-03). Cross terms with constant mean metric or shear integrate to zero by periodicity. This gives the three momentum conditions in (W-03). The homogeneous linear equations keep $h,\Sigma$ constant, while the wave equation and periodic integration conserve wave energy and momentum.

The integrated conditions are the compact linearization-stability conditions \cite{BrillDeser1973,ArmsMarsdenMoncrief1982}. Positive TT energy excludes pure-wave directions but permits mixed directions with volume velocity. We next state their sufficiency on the fixed flat torus and give explicit second-order corrections.


\subsection{Application of the known theorem on the fixed flat torus}

\stepcounter{equation}\begin{equation}\tag{WD-04}\label{eq:WD-04}
\partial_i\partial_j\nu-\delta_{ij}\Delta\nu=0,\qquad
\partial_iX_j+\partial_jX_i=0,\qquad
\det(\mathcal S_k^\dagger\mathcal S_k)=32(k^2)^5\quad(k\ne0).
\end{equation}

In units $q_0=1$, (WD-04) describes the adjoint kernel. Taking the trace and using periodicity makes $\nu$ constant; the periodic Killing equation makes $X$ constant. These exhaust the four KIDs. Coordinate-linear rotations and boosts do not descend to the fixed lattice. The last identity checks the combined four-column symbol at nonzero wavevector.

\stepcounter{equation}\begin{equation}\tag{WD-05}\label{eq:WD-05}
\begin{aligned}
\delta g&=C+2t(hI+\Sigma)+\gamma, &\delta\pi&=\Sigma-2hI+\tfrac12\dot\gamma,\\
\beta&=\mathrm{tr}(\delta\pi)=-6h,&\omega^{\rm TT}&=\Sigma+\tfrac12\dot\gamma,\\
\langle(\omega^{\rm TT})^2\rangle-\beta^2/6
+\tfrac14\langle(\partial\gamma)^2\rangle
&=\mathrm{tr}\Sigma^2-6h^2+\tfrac14\langle\dot\gamma^2+(\partial\gamma)^2\rangle.
\end{aligned}
\end{equation}

$C$ and $\Sigma$ are constant, $\Sigma$ is tracefree and $\gamma$ is mean-zero TT. Equation (WD-05) identifies the Hamiltonian cone, Eq. (3.5) of Arms--Marsden--Moncrief, with the energy condition in (W-03); their momentum cones, Eq. (3.1), agree with (WD-03). Their Section 3 treats exactly a flat $g_0$ with $\pi_0=0$, and their introduction includes the flat spacetime $T^3\times\mathbb R$ with four Killing fields. We apply their Theorem 3.3 (Section 3, p. 96) to smooth periodic data on a compact boundaryless CMC slice. Curves may be interpreted, conservatively, in $H^s\times H^{s-1}$ for integer $s\ge5$, without claiming minimal regularity. The four conditions then characterize tangents to exact constraint-data curves. Local Cauchy development and the LC lift give EC vacuum families \cite{ArmsMarsdenMoncrief1982}.

This concerns a local initial-value family on a common time interval, not a two-time boundary-value problem, all-time regularity or an exact family closed in a finite Fourier space. Higher-order corrections may contain new harmonics.

\subsection{Explicit second-order corrections and two-wave closure}

For the standing wave (W-04), expanding each exponent as $\epsilon A_1+\epsilon^2A_2$, and similarly for $B,P$, makes the first and second coefficients of all ten Ricci components vanish, for either sign of $h$. This is a local substitution check through second order, not an exact finite-amplitude solution.

\stepcounter{equation}\begin{equation}\tag{WD-06}\label{eq:WD-06}
\begin{aligned}
f&=a\cos x\cos t,\quad g=b\cos y\cos t,\quad h^2=(a^2+b^2)/24,\\
H_{\mu\nu}&=\mathrm{diag}(-6ht,\,2ht+g,\,2ht+f,\,2ht-f-g),\\
\mathcal K&=\{(0,0),(\pm1,0),(0,\pm1),(\pm2,0),(0,\pm2),(\pm1,\pm1)\}.
\end{aligned}
\end{equation}

\stepcounter{equation}\begin{equation}\tag{WD-07}\label{eq:WD-07}
\begin{aligned}
\ddot{\bar Q}_{\mu\nu,n}+|n|^2\bar Q_{\mu\nu,n}&=-2G^{[2]}_{\mu\nu,n}[H],\\
-\dot{\bar Q}_{0\nu}+\partial_i\bar Q_{i\nu}&=0,\\
\psi&=-\tfrac7{16}(a^2\cos2x+b^2\cos2y)-ab\cos x\cos y,\\
\Delta\psi&=-2G^{[2]}_{00}[H](0),\\
\bar Q_{ij}(0)&=\delta_{ij}\psi,\qquad \bar Q_{00}(0)=\bar Q_{0i}(0)=0,\\
\dot{\bar Q}_{ij}(0)=\dot{\bar Q}_{00}(0)&=0,\qquad \dot{\bar Q}_{0i}(0)=\partial_i\psi.
\end{aligned}
\end{equation}

Write $g=\eta+\epsilon H+\epsilon^2Q$, with bars denoting trace reversal by $\eta$. Source conservation and the initial conditions in (WD-07) preserve the second-order harmonic condition. Keeping all thirteen wavevectors, integration of the forced oscillators constructs $Q$ satisfying the local Einstein equations and reality through second order. The public \texttt{two\_waves.json} supplies the source and correction coefficients. Removing the zero wavevector does not give the same construction.
